% =====================================================================
\phantomsection
\addcontentsline{toc}{part}{Languages for Technical and Scientific Computing}
\BookChapter[MATLAB and GNU Octave]{MATLAB and GNU Octave: The Matrix Laboratory}
\label{ch:matlab}

\begin{researchquestion}
What changes when every variable is a matrix?
\end{researchquestion}
\OrientationNote
% =====================================================================

\begin{factsheet}{MATLAB}
First appeared & late 1970s (Fortran version); commercial MATLAB 1984; two releases per year (named, for example, R2026a and R2026b)\\
Designed by & Cleve Moler (University of New Mexico); company MathWorks founded 1984 with Jack Little and Steve Bangert\\
Paradigms & array programming, procedural, object-oriented\\
Typing & dynamic; everything is a matrix (default type \code{double})\\
Execution & interpreted with a JIT-accelerated execution engine\\
File extensions & \file{.m} (scripts, functions), \file{.mlx} (live scripts), \file{.mat} (data), \file{.slx} (Simulink)\\
Cost & commercial; campus-wide licences common at universities. Free alternative: GNU Octave\\
Used for & control systems, signal and image processing, numerical analysis, data analysis,
  system modelling and simulation (Simulink), automotive and aerospace embedded code generation, teaching\\
Famous software & Simulink, Control System and Signal Processing Toolboxes, Stateflow; Simulink models
  in most automotive and aerospace companies\\
\end{factsheet}

MATLAB, short for \emph{MATrix LABoratory}, has been the standard computational tool of
engineering education and of many engineering industries for four decades. Its central
idea is simple and powerful: the basic data type is the matrix, and the language is
designed to make linear algebra and array operations read like mathematics. If you
continue in control, signal processing, mechatronics, automotive or aerospace engineering,
you will very likely use MATLAB and Simulink.

\LearningOutcomes{Create, index and combine vectors and matrices; choose correctly between matrix and element-wise operators; write scripts, function files and anonymous functions; solve linear systems and ordinary differential equations with built-in tools; describe what Simulink and the toolboxes are for; and translate simple MATLAB code into Python.}

\section{History and philosophy}

\begin{history}
In the 1970s Cleve Moler, a numerical analyst, helped write the Fortran libraries LINPACK
(linear equations) and EISPACK (eigenvalues). He wanted his students at the University of
New Mexico to use these libraries without having to learn Fortran, so he wrote a small
interactive program, MATLAB, in which you could type \code{A = [1 2; 3 4]} and \code{eig(A)}.
Copies spread to other universities. In 1983 the engineer Jack Little saw its commercial
potential, rewrote it in C with Steve Bangert, and in 1984 they founded The MathWorks with
Moler. Control engineers adopted it quickly, and toolboxes for control design and signal
processing followed. Simulink, a graphical block-diagram environment for simulating
dynamic systems, was added in 1990.
\end{history}

MATLAB's philosophy is ``the matrix is the computer''. Every variable is an array; a
scalar is a $1\times1$ matrix. Operators such as \code{*} and \code{\textbackslash} have
their linear-algebra meaning. The language is complemented by a polished environment
(editor, debugger, workspace browser, plotting) and by a large collection of professionally
developed and documented \term{toolboxes}. You pay for this quality: MATLAB is commercial
software and each toolbox has its own licence.

\textbf{GNU Octave} is a free, open-source program that implements most of the core MATLAB
language. All examples in this chapter were tested in Octave; they also run in MATLAB.

\section{Getting started}

The MATLAB desktop has a \emph{Command Window} for interactive work, an \emph{Editor} for
scripts and functions, a \emph{Workspace} panel showing current variables, and a
\emph{Current Folder} panel. MATLAB Online runs in a web browser.
\begin{lstlisting}[language=Matlab]
>> 2 + 3
ans =
     5
>> g = 9.81;              % a semicolon suppresses the output
>> m = 80
m =
    80
>> sqrt(m*g/0.25)
ans =
   56.0286
\end{lstlisting}
Note that a statement \emph{without} a semicolon prints its result. This is handy
interactively but floods the screen in scripts, so most lines in a script end with
\code{;}. The variable \code{ans} holds the last unnamed result. Comments start with
\code{\%}.

A \term{script} is a file of commands, run by typing its name (without \file{.m}) or
pressing Run:
\codefile{Matlab}{matlab/hello.m}{a first script}
\outfile{matlab/hello.out}

\section{Everything is a matrix}

\subsection*{Build arrays with brackets, colons and helper functions}

\begin{lstlisting}[language=Matlab]
x = [1 2 3]             % row vector (1 x 3); commas also work: [1, 2, 3]
y = [1; 2; 3]           % column vector (3 x 1)
A = [1 2; 3 4]          % 2 x 2 matrix: ; starts a new row
t = 0:0.5:2             % 0  0.5  1.0  1.5  2.0  (end INCLUDED)
u = linspace(0, 1, 5)   % 5 equally spaced points from 0 to 1
Z = zeros(3)            % 3 x 3 zeros;  zeros(2, 4) is 2 x 4
O = ones(2, 3); I = eye(3); R = rand(2);  M = magic(4);
size(A), numel(t), length(t)
\end{lstlisting}

\subsection*{Indices start at 1 and use round brackets}

Indices start at \textbf{1} and use round brackets. The keyword \code{end} means the last
index, and a colon alone means ``all'':
\begin{lstlisting}[language=Matlab]
t(1)          % first element
t(end)        % last element
t(2:4)        % elements 2, 3, 4
A(2, 1)       % row 2, column 1
A(:, 2)       % whole second column
A(1, :)       % whole first row
t(t > 1)      % logical indexing: all elements greater than 1
\end{lstlisting}
Arrays grow automatically if you assign beyond their end (\code{t(10) = 7} pads with
zeros), which is convenient but slow inside loops. \emph{Preallocate} with \code{zeros}
when you know the final size.

\subsection*{A dot decides between matrix and element-wise operations}

This is the single most important thing to understand about MATLAB syntax.
\begin{center}
\small
\begin{tabular}{@{}lll@{}}
\toprule
\textbf{Operation} & \textbf{Matrix (linear algebra)} & \textbf{Element by element}\\
\midrule
multiplication & \code{A*B} & \code{A.*B}\\
division & \code{A/B}, \code{A\textbackslash B} & \code{A./B}\\
power & \code{A\^{}2} (= \code{A*A}) & \code{A.\^{}2}\\
transpose & \code{A\sq{}} (conjugate), \code{A.\sq{}} & ---\\
\bottomrule
\end{tabular}
\end{center}
Addition and subtraction are always element-wise. If you want to square each element of a
vector of times, you must write \code{t.\^{}2}; \code{t\^{}2} tries to compute the matrix
product \code{t*t} and fails because a $1\times5$ matrix cannot be multiplied by itself.

\begin{pitfall}[Common pitfall: the missing dot]
The error message ``Incorrect dimensions for matrix multiplication'' almost always means
you forgot a dot: \code{x*y} instead of \code{x.*y}, or \code{1/x} instead of
\code{1./x}. Conversely, if you \emph{do} want a matrix product and get an unexpected
result, check that you did not add a dot by habit.
\end{pitfall}

\codefile{Matlab}{matlab/matrices.m}{matrices and linear algebra}
\outfile[title={\footnotesize\sffamily\color{muted}Output (Octave)}]{matlab/matrices.out}
The backslash operator \code{A\textbackslash b} solves $A\mathbf{x} = \mathbf{b}$. It is
one of MATLAB's cleverest features: it examines the matrix and automatically chooses an
appropriate LAPACK algorithm (Cholesky for symmetric positive-definite matrices, LU for
general square matrices, QR for least-squares problems with rectangular matrices).

\section{Control flow}

Blocks are closed with \code{end}. Indentation is optional but recommended.
\begin{lstlisting}[language=Matlab]
if Re < 2300
    regime = 'laminar';
elseif Re < 4000
    regime = 'transitional';
else
    regime = 'turbulent';
end

for k = 1:10              % k = 1, 2, ..., 10
    s = s + 1/k^2;
end

while abs(dx) > 1e-12
    dx = f(x)/df(x);
    x = x - dx;
end

switch material
    case 'steel'
        E = 210e9;
    case {'aluminium', 'aluminum'}
        E = 69e9;
    otherwise
        error('Unknown material %s', material);
end
\end{lstlisting}
Logical operators are \code{\&\&} and \code{||} for scalars, \code{\&} and \code{|} for
element-wise operations on arrays, and \code{\textasciitilde} for ``not''; ``not equal'' is
\code{\textasciitilde=}.

\section{Functions}

A function lives in its own file, whose name must match the function name:
\codefile{Matlab}{matlab/terminal_velocity.m}{a function file}
The comment lines directly after the \code{function} line are the help text, shown by
\code{help terminal\_velocity}. \code{nargin} is the number of arguments actually passed,
which is MATLAB's way of providing default values. Because the function uses element-wise
operators (\code{.*}, \code{./}), it works for arrays as well as scalars:
\codefile{Matlab}{matlab/use_function.m}{calling the function}
\outfile{matlab/use_function.out}
Note how \code{fprintf} recycles its format string for each column of the matrix
\code{[masses; vt]}: this is a compact MATLAB idiom for printing tables. Functions can return
several outputs, \code{[mu, sigma] = stats(x)}, and modern MATLAB also allows
\term{local functions} at the end of a script file and an \code{arguments} block for
validating inputs.

Short functions can be written inline as \term{anonymous functions}:
\begin{lstlisting}[language=Matlab]
f  = @(x) x.^3 - 2*x - 5;
df = @(x) 3*x.^2 - 2;
root = fzero(f, 2)          % built-in root finder: 2.0946
area = integral(@sin, 0, pi)% built-in integration: 2.0000
\end{lstlisting}
The \code{@} creates a \term{function handle}, which can be passed to other functions such
as \code{fzero}, \code{integral}, \code{fminsearch} and the ODE solvers.

\section{Case study: the falling skydiver}

\codefile{Matlab}{matlab/drag.m}{Euler solution of the drag equation}
\outfile{matlab/drag.out}
Compare this with the \py version. The time vector \code{0:dt:t\_end} includes the end
point; the loop starts at index 1; and because the anonymous function \code{exact\_speed}
uses \code{.\^{}} and works element-wise, the exact solution at all 41 times is computed in
one line. The \code{fprintf} statement prints the whole table at once from a 4-row matrix.

\section{Plotting}

Plotting is one of MATLAB's traditional strengths, and Matplotlib's interface was modelled
on it:
\begin{lstlisting}[language=Matlab]
t = linspace(0, 20, 400);
vt = sqrt(80*9.81/0.25);
plot(t, vt*tanh(9.81*t/vt), 'b-', 'LineWidth', 1.5)
hold on                                  % add to the same axes
yline(vt, '--', 'terminal velocity');
hold off
xlabel('t (s)'); ylabel('v (m/s)');
title('Falling skydiver'); grid on; box on
legend('exact', 'Location', 'southeast')
exportgraphics(gcf, 'drag.pdf')          % vector graphics for reports
\end{lstlisting}
Other common plot types are \code{semilogy}, \code{loglog}, \code{bar}, \code{histogram},
\code{scatter}, \code{subplot} and \code{tiledlayout} (several plots in one figure),
\code{surf}, \code{mesh} and \code{contour} (surfaces and fields).

\section{Solving ODEs}

MATLAB's ODE solvers are widely used in engineering courses. \code{ode45} is an adaptive
Runge--Kutta method (the counterpart of SciPy's \code{solve\_ivp} with its default
method); \code{ode15s} is for \emph{stiff} problems such as chemical kinetics. A system of
ODEs is passed as a function handle returning a column vector of derivatives. Here is a
damped mass--spring system $m\ddot x + b\dot x + kx = 0$, written as two first-order
equations for $x$ and $\dot x$:
\begin{lstlisting}[language=Matlab]
m = 1.0; b = 0.4; k = 4.0;
rhs = @(t, y) [y(2); -(b*y(2) + k*y(1))/m];   % y = [x; xdot]
[t, y] = ode45(rhs, [0 20], [1; 0]);           % from x = 1 at rest
plot(t, y(:,1)); xlabel('t (s)'); ylabel('x (m)');
\end{lstlisting}

\section{Toolboxes and Simulink}

MATLAB's toolboxes package decades of engineering methods: Control System (transfer
functions, Bode plots, controller design), Signal Processing (filters, spectra),
Optimization, Statistics and Machine Learning, Image Processing, Partial Differential
Equations, Curve Fitting, Symbolic Math, and many domain-specific ones (Aerospace,
Powertrain, Robotics, Communications). For example, a control engineer can write
\begin{lstlisting}[language=Matlab]
s = tf('s');
G = 1/(s^2 + 0.4*s + 4);     % plant transfer function
C = pid(2, 1, 0.5);          % PID controller
T = feedback(C*G, 1);        % closed loop
step(T)                      % plot the step response
\end{lstlisting}

\textbf{Simulink} is a graphical environment in which a dynamic system is drawn as a block
diagram --- integrators, gains, sums, transfer functions, lookup tables, switches --- and
simulated. It is used throughout the automotive and aerospace industries for
\term{model-based design}: engineers build and test a controller as a Simulink model, and
code generators (Embedded Coder) then produce C code from the model that runs on the
vehicle's electronic control unit. A large share of the software in modern cars is
generated this way. Simscape extends Simulink with physical components (mechanical,
electrical, hydraulic, thermal), using an equation-based approach similar to Modelica
(Section~\ref{sec:modelica}).

\section{MATLAB, Octave and Python}

Many engineering students learn both MATLAB and \py, and the two are closely related.
Table~\ref{tab:matlab-python} translates common operations. The biggest differences are
indexing (1 versus 0), slices (end included versus excluded), and the meaning of
\code{*}.

\begin{table}[htbp]
\centering
\caption{MATLAB and \py/NumPy side by side.}
\label{tab:matlab-python}
\small
\begin{tabular}{@{}ll@{}}
\toprule
\textbf{MATLAB} & \textbf{\py with NumPy} (\code{import numpy as np})\\
\midrule
\code{a = [1 2 3];} & \code{a = np.array([1, 2, 3])}\\
\code{A = [1 2; 3 4];} & \code{A = np.array([[1, 2], [3, 4]])}\\
\code{a(1)}, \code{a(end)} & \code{a[0]}, \code{a[-1]}\\
\code{A(2, :)} & \code{A[1, :]}\\
\code{a(2:3)} & \code{a[1:3]}\\
\code{0:0.5:2} & \code{np.arange(0, 2.5, 0.5)} or \code{np.linspace(0, 2, 5)}\\
\code{zeros(3, 2)} & \code{np.zeros((3, 2))}\\
\code{A * B} & \code{A @ B}\\
\code{A .* B} & \code{A * B}\\
\code{A .\^{} 2} & \code{A**2}\\
\code{A\sq{}} & \code{A.T} (or \code{A.conj().T})\\
\code{A \textbackslash\ b} & \code{np.linalg.solve(A, b)}\\
\code{eig(A)} & \code{np.linalg.eig(A)}\\
\code{size(A)}, \code{numel(A)} & \code{A.shape}, \code{A.size}\\
\code{sum(A, 1)}, \code{sum(A, 2)} & \code{A.sum(axis=0)}, \code{A.sum(axis=1)}\\
\code{f = @(x) x.\^{}2;} & \code{f = lambda x: x**2}\\
\code{fzero(f, x0)} & \code{scipy.optimize.fsolve(f, x0)} / \code{brentq}\\
\code{ode45(f, [t0 t1], y0)} & \code{scipy.integrate.solve\_ivp(f, (t0, t1), y0)}\\
\code{plot(x, y)} & \code{plt.plot(x, y)}\\
\code{load data.mat} & \code{scipy.io.loadmat("data.mat")}\\
\bottomrule
\end{tabular}
\end{table}

\section{Strengths and weaknesses}

\begin{center}
\small
\begin{tabularx}{\linewidth}{@{}>{\raggedright\arraybackslash}X>{\raggedright\arraybackslash}X@{}}
\toprule
\textbf{Strengths} & \textbf{Weaknesses}\\
\midrule
Natural notation for matrices and linear algebra & Expensive licences; code cannot be shared with people who lack them\\
Excellent integrated environment and documentation & Weak as a general-purpose language (strings, data structures, web)\\
Professional toolboxes; Simulink for model-based design & One-function-per-file rule makes large projects awkward\\
Standard in control, signal processing, automotive and aerospace & Loops historically slow (much improved by the JIT)\\
Code generation for embedded targets & Closed source: you cannot inspect how built-in functions work\\
\bottomrule
\end{tabularx}
\end{center}

\section{Further reading}
\begin{itemize}
\item MathWorks, \emph{MATLAB Onramp}, a free 2-hour online tutorial,
  \url{https://matlabacademy.mathworks.com}.
\item C.~Moler, \emph{Numerical Computing with MATLAB}, SIAM, 2004 (free online).
\item S.~J.~Chapman, \emph{MATLAB Programming for Engineers}, 6th ed., Cengage, 2019.
\item GNU Octave documentation, \url{https://octave.org/doc/}.
\end{itemize}

\begin{exercises}
\item Install GNU Octave (or use MATLAB Online) and run \file{drag.m}. Then add the plot
  commands at the end of the file and produce a figure similar to
  Figure~\ref{fig:dragresult}.
\item Given \code{x = 0:0.1:1}, which of the following produce errors and why?
  \code{x*x}, \code{x.*x}, \code{x*x\sq{}}, \code{x\sq{}*x}, \code{x\^{}2}, \code{2.\^{}x}.
  What are the sizes of the results that do work?
\item Write a function \code{[mu, s] = mystats(x)} that returns the mean and the sample
  standard deviation of a vector without using \code{mean} or \code{std}. Test it against the
  built-in functions with the data in \file{yield.txt} (\code{x = load(\sq{}../data/yield.txt\sq{})}).
\item Use \code{ode45} to solve the drag equation and compare with the Euler result at
  $t = 10\unit{s}$.
\item Solve the beam problem from Chapter~\ref{ch:fortran} (cantilever deflection) in
  MATLAB with a single vectorised expression, and plot $w(x)$.
\item Translate \file{matrices.m} into \py with NumPy, using
  Table~\ref{tab:matlab-python}.
\end{exercises}
